\section{Research question}
\label{sec:question}

None of the six papers reviewed here tests provenance verification, independent rebuilding, and DNS monitoring together. The proof-of-concept (PoC) uses DNS monitoring for detection only; the resolver block is an operator action, not a monitoring-driven defence. It compares an unsigned catalogue digest and verifies the package signature, but it implements no independent verification: no rebuild and no reference-artefact comparison. The article therefore asks:

\begin{quote}
\textbf{Which stages of a compromised but legitimately signed software-update workflow can independent build-artefact verification and DNS activity monitoring identify in a controlled Mininet and Docker lab, and how many false positives does each control produce?}
\end{quote}

Three sub-questions separate the stages. Before release, can an independently generated reference artefact or a trusted provenance policy distinguish the malicious V3 from an approved release when both pass signature checks? After installation, which check-ins and CNAME handoffs does DNS logging reveal, and how many alerts does a stated rule raise on benign runs? And which stages remain invisible when either control is used alone?

The evaluation would run a baseline approved update, the tampered update, an artefact comparison, DNS monitoring on labelled runs, and both controls together. It would record detected releases, flagged check-ins, false positives on benign runs, and the stage at which each alert fires. The lab has no benign DNS workload at present, so the false-positive half of the question is not yet measurable; benign comparison traffic is required before any detection rate means anything. If reproducible builds cannot be implemented, an approved-reference comparison is the narrower fallback, and it must not be described as proof of reproducible-build security.

The limits are the lab's: a predefined check-in routine, simplified DNS, no SUNSPOT replacement, no original DGA, and no later-stage malware. Results would describe this environment, not SolarWinds.
